\section{Consonantes especiales (g, ch, sch, r)}
\textbf{Qué son:} Consonantes con fricación dorsal o combinaciones típicas del neerlandés.

\textbf{Por qué importan:} Distinguen acento y comprensión; \textit{g/ch} es un sonido clave.

\textbf{Cómo practicarlas:}
\begin{itemize}[leftmargin=*]
  \item \textit{g/ch}: fricativa velar/uvular sorda. Palabras: \textit{goed, graag, licht, lachen}.
  \item \textit{sch}: se pronuncia /sx/ al inicio: \textit{school, schip}. En medio puede sonar /s/ + /x/ suave.
  \item \textit{r}: puede ser uvular o alveolar; elige una y sé consistente. Palabras: \textit{rood, straat, leren}.
  \item \textit{v/z}: sonoras suaves; cuida que no suenen /b/s.
\end{itemize}

\textbf{Ejercicios:}
\begin{itemize}[leftmargin=*]
  \item Soplo de aire: practica \textit{g} exhalando al empañar un espejo mientras dices \textit{goed}.
  \item Trenes de palabras: \textit{graag-gaan-groen}; \textit{school-schip-schrift}; \textit{rood-reizen-leren}.
\end{itemize}
